\scanpage{6}

\uline{Icosaèdre}

I Sa description concrète et ses premières propriétés

1. Description \inserted{grossière} par un modèle concret (ou dessin)

\deleted{2.} Énumération

\diagram{$\begin{array}{lll}\text{Nb de sommets}&s=12&\\
\text{d'arêtes}&a=30&2a=ps=qf\\
\text{faces}&f=20&p=5,\ q=3\end{array}$}

2. Description des propriétés polyédrales

\margin{parenthèse : polyèdre combinatoire}
% slanted

A) Description comme ``géométrie d'incidences''

B) \rule{3cm}{0.35pt} ``polyèdre combinatoire'' \inserted{par \deleted{propriétés} \inserted{existence} des $\operatorname{Sup}$ \inserted{majorés} et $\operatorname{Inf}$ min.}

3. Distances combinatoires \unsure{,} \deleted{\unsure{aux}} \inserted{de sommets}, antipodisme de \inserted{sommets}

4. Perspectives

A) Par sommet ou couple sommets opposés (gpe d'autom. $\mathbb D_5$ ou $\mathbb Z_2\times\mathbb D_5$)

\margin{\unsure{on trouve|introduire|Interprétation} l'antipodisme comme \uline{automorphisme combinatoire}}
% slanted

\uline{Parenthèse} groupe des automorphismes du polygone combinatoire ....

Transitivité du groupe des automorphismes sur sommets, et $=$ sur couples $(s,a)$ $(s<a)$ et $=$ sur triples $(s<a<f)$

On dit que l'icosaèdre est un polyèdre combinatoire régulier (\uline{NB} Il y a une vaste infinité d'autres polyèdres combinatoires réguliers -- mais parmi ceux-ci, cinq seulement : resp. \unsure{peuv$^{t}$|peut} se ``réaliser'' comme des polyèdres convexes dans l'espace euclidien de dim 3 ....

% archive sheet number: 5

